% 
\documentclass[twoside, 11pt, a4paper]{article}

\usepackage[a4paper, top=2.5cm, bottom=2.5cm, left=2cm, right=2cm]{geometry}
\usepackage[english]{babel}

\setlength{\parindent}{0 in}
\setlength{\parskip}{0.1 in}

%
% ADD PACKAGES here:
%
\usepackage{amsmath,amsfonts,graphicx}

\newcounter{lecnum}
\renewcommand{\thepage}{\thelecnum-\arabic{page}}
\renewcommand{\thesection}{\thelecnum.\arabic{section}}
\renewcommand{\theequation}{\thelecnum.\arabic{equation}}
\renewcommand{\thefigure}{\thelecnum.\arabic{figure}}
\renewcommand{\thetable}{\thelecnum.\arabic{table}}

%
% The following macro is used to generate the header.
%
\newcommand{\lecture}[2]{
   \pagestyle{myheadings}
   \thispagestyle{plain}
   \newpage
   \setcounter{lecnum}{#1}
   \setcounter{page}{1}
   \noindent
   \begin{center}
   \framebox{
      \vbox{\vspace{2mm}
    \hbox to 6.28in { {\bf EE502 - Linear Systems Theory II
	\hfill Spring 2023} }
       \vspace{4mm}
       \hbox to 6.28in { {\Large \hfill Lecture #1 \hfill} }
       \vspace{2mm}
       \hbox to 6.28in { {\it Lecturer: #2 \hfill } }
      \vspace{2mm}}
   }
   \end{center}
   \markboth{Lecture #1}{Lecture #1}

   \vspace*{4mm}
}

\renewcommand{\cite}[1]{[#1]}
\def\beginrefs{\begin{list}%
        {[\arabic{equation}]}{\usecounter{equation}
         \setlength{\leftmargin}{2.0truecm}\setlength{\labelsep}{0.4truecm}%
         \setlength{\labelwidth}{1.6truecm}}}
\def\endrefs{\end{list}}
\def\bibentry#1{\item[\hbox{[#1]}]}


\newcommand{\fig}[3]{
			\vspace{#2}
			\begin{center}
			Figure \thelecnum.#1:~#3
			\end{center}
	}

% Use these for theorems, lemmas, proofs, etc.
\newtheorem{theorem}{Theorem}[lecnum]
\newtheorem{lemma}[theorem]{Lemma}
\newtheorem{proposition}[theorem]{Proposition}
\newtheorem{claim}[theorem]{Claim}
\newtheorem{corollary}[theorem]{Corollary}
\newtheorem{definition}[theorem]{Definition}
\newenvironment{proof}{{\bf Proof:}}{\hfill\rule{2mm}{2mm}}
\newtheorem{exmp}[theorem]{Ex}

% **** IF YOU WANT TO DEFINE ADDITIONAL MACROS FOR YOURSELF, PUT THEM HERE:

\begin{document}

% Lecture Details
\lecture{13}{Asst. Prof. M. Mert Ankarali}

%%%%%%%%%%%%%%%%%%%%%%%%%%

\section{State Feedback}

\textbf{Intuition \& Insight:} By feeding back the state variables into the input, we are fundamentally altering the system's internal vector field. The primary goal of state feedback is to reshape the transient response of the system (e.g., settling time, overshoot, damping) and ensure asymptotic stability by strategically relocating the closed-loop poles (eigenvalues). 

State-feedback-based control policies for LTI systems start with the assumption that 
we have ``access'' to all of the states of the system, either via direct measurement 
or through some observer/estimator/tracker. In that context, a family of state-feedback 
controllers for CT and DT LTI systems can be constructed as:
%
\begin{align*}
 u(t) = \gamma r(t) - K x(t) \quad \& \quad   u[k] = \gamma r[k] - K x[k]
\end{align*}
%
where $r(t)$ can be considered as the reference signal (most of the time it is), 
$\gamma$ is a feed-forward scaling factor, and $K$ is the state-feedback gain matrix. 
Now let's find a state-space representation for the dynamics of the closed-loop system
for both CT and DT LTI systems under the state-feedback rule proposed above:
%
\begin{align*}
	\dot{x} &= A x + B \left( \gamma r(t) - K x(t) \right) \implies \dot{x} = \left( A - B K \right) x + \gamma B r(t)
	\\
	x[k+1] &= A x[k] + B \left( \gamma r[k] - K x[k] \right) \implies x[k+1] = \left( A - B K \right) x[k] + \gamma B r[k]
\end{align*}
%
In both cases, the closed-loop system matrix and input matrix take the following form:
%
\begin{align*}
	A_c = A - B K \ , \  B_c = \gamma B
\end{align*}
%
A key theoretical question in this domain is whether we can find a gain matrix $K$ such that the eigenvalues of $A_c$ are located at \textit{arbitrary} desired locations in the complex plane (subject to complex conjugate pairing). 

\textbf{Theorem: (Eigenvalue/Pole Placement)} Given the pair $(A,B)$, there exists a feedback gain $K$ such that the closed-loop characteristic polynomial matches any arbitrarily chosen desired polynomial:
%
\begin{align*}
	\mathrm{det}\left[ \lambda I - (A - B K ) \right] &= \lambda^n + a_{n-1}^* \lambda^{n-1}
	+ \cdots + a_{1}^* \lambda + a_0^*
	\\
	\forall \mathcal{A}^* &= \lbrace a_0^* , \, a_1^* \, \dots \, a_{n-1}^* \rbrace \subset \mathbb{R}
\end{align*} 
%
if and only if $(A,B)$ is reachable. 

\textbf{Proof:} For a general complete proof, we need to show that the reachability of $(A,B)$ is both necessary and sufficient. 

\textbf{Proof of Necessity:} Let's assume that $(A,B)$ is \textit{not} reachable. By the PBH test for reachability, there exists a left eigenvector $w_u^T \neq 0$ associated with an unreachable eigenvalue $\lambda_u$ such that
$w_u^T A = \lambda_u w_u^T $ and $w_u^T B = 0$. 

Now let's check whether $w_u^T$ remains a left eigenvector for the closed-loop matrix $A_c$:
%
\begin{align*}
	w_u^T A_c &= w_u^T (A - B K ) = w_u^T A - w_u^T B K = \lambda_u w_u^T - 0 \cdot K = \lambda_u w_u^T
	\\
	w_u^T B_c &= w_u^T B \gamma = 0
\end{align*} 
%
Here, not only have we shown that the unreachable mode $\lambda_u$ cannot be moved (hence it is impossible to place the poles at arbitrary locations, proving necessity), but we have also shown that state-feedback rules do not alter the system's fundamental reachability properties.

\textbf{Proof of Sufficiency:} We will constructively prove sufficiency for a single-input single-output (SISO) case, i.e., $B \in \mathbb{R}^{n \times 1}$; however, the reader should note that for a complete general proof, the multi-input case must be analyzed (typically utilizing Luenberger controllable canonical forms). 

Let's assume that $(A,B)$ is reachable. We know that reachability is invariant under similarity transformations. Let's apply a transformation $z = T^{-1} x$ (where $\mathrm{det}(T) \neq 0$), yielding:
%
\begin{align*}
	\dot{z} = \bar{A} z + \bar{B} u \quad \text{where} \quad \bar{A} = T^{-1} A T \ , \ \bar{B} = T^{-1} B 
\end{align*}
%
Since $(A,B)$ is reachable, its reachability matrix $\mathbf{R}$ is full rank. Let's intentionally choose our transformation matrix $T$ to be this reachability matrix:
%
\begin{align*}
T = \mathbf{R} = \left[ \begin{array}{c|c|c|c|c} A^{n-1} B & A^{n-2} B & \cdots & A B &  B \end{array} \right]
\end{align*}
%
Then the transformed input matrix $\bar{B}$ can be isolated:
% 
\begin{align*}
B = T \bar{B} = \mathbf{R} \bar{B} \implies \bar{B} = \mathbf{R}^{-1} B = \begin{bmatrix} 0 \\ 0 \\ \vdots \\ 0 \\ 1 \end{bmatrix}
\end{align*}
%
(Since $B$ is exactly the last column of $\mathbf{R}$, multiplying $\mathbf{R}^{-1}$ by $B$ plucks out the last canonical basis vector). Similarly, $\bar{A}$ can be expressed by analyzing $A \mathbf{R} = \mathbf{R} \bar{A}$:
\begin{align*}
A \left[ \begin{array}{c|c|c|c|c} A^{n-1} B & A^{n-2} B & \cdots & A B &  B \end{array} \right] &= \left[ \begin{array}{c|c|c|c|c} A^{n-1} B & A^{n-2} B & \cdots & A B &  B \end{array} \right] \bar{A}
\\
\left[ \begin{array}{c|c|c|c|c} A^{n} B & A^{n-1} B & \cdots & A B &  B \end{array} \right] &= \mathbf{R} \left[ \begin{array}{c|c|c|c|c|c} 
\bar{a}_{11} & 1 & 0 & 0 & \cdots & 0 \\
\bar{a}_{12} & 0 & 1 &  0 & \cdots & 0 \\
\vdots & \vdots & \vdots & \ddots &  & \vdots \\
\bar{a}_{1(n-2)} & 0 & 0 & \cdots  & 1 & 0 \\
\bar{a}_{1(n-1)} & 0 & 0 & \cdots  & 0 & 1 \\
\bar{a}_{1(n)} & 0 & 0 & \cdots  & 0 & 0 \\
\end{array} 
\right] 
\end{align*}
%
The first column coefficients $\bar{a}_{1i}$ are found by expressing $A^n B$ as a linear combination of the lower powers using the Cayley-Hamilton theorem:
%
\begin{align*}
A^{n} &= - \left( a_{n-1} A^{n-1} + a_{n-2} A^{n-2} + \cdots + a_1 A + a_0 I \right) 
\end{align*}
where $a_i$ are the coefficients of the open-loop characteristic polynomial:
\begin{align*}
\mathrm{det}[\lambda I - A] &= \lambda^n + a_{n-1}  \lambda^{n-1} + \cdots + a_{1} \lambda + a_0
\end{align*}
Substituting this, $\bar{A}$ takes a block companion form:
\begin{align*}
\bar{A} &= \left[ \begin{array}{c|c|c|c|c|c} 
-a_{n-1} & 1 & 0 & 0 & \cdots & 0 \\
-a_{n-2} & 0 & 1 &  0 & \cdots & 0 \\
\vdots & \vdots & \vdots & \ddots &  & \vdots \\
-a_{2} & 0 & 0 & \cdots  & 1 & 0 \\
-a_{1} & 0 & 0 & \cdots  & 0 & 1 \\
-a_{0} & 0 & 0 & \cdots  & 0 & 0 
\end{array} 
\right] 
\end{align*}
%
\textbf{Transformation to Reachable Canonical Form: A Graphical Intuition}

Note that the transformation $T = \mathbf{R}$ yielded a specific companion form $(\bar{A}, \bar{B})$. While mathematically valid, we want to transform this into the (more useful) \textbf{reachable canonical form} $(\tilde{A}, \tilde{B})$, which cleanly exposes the characteristic polynomial coefficients entirely in the \textit{last row}:
%
\begin{align*}
\tilde{A} = \left[ \begin{array}{c|c|c|c|c|c} 
0 & 1 & 0 & 0 & \cdots & 0 \\
0 & 0 & 1 &  0 & \cdots & 0 \\
\vdots & \vdots & \vdots & \ddots &  & \vdots \\
0 & 0 & 0 & \cdots  & 1 & 0 \\
0 & 0 & 0 & \cdots  & 0 & 1 \\
-a_{0} & -a_{1} & -a_{2} & \cdots  & -a_{n-2} & -a_{n-1} 
\end{array} 
\right] 
\ , \ \tilde{B} = \begin{bmatrix} 0 \\ 0 \\ \vdots \\ 0 \\ 1 \end{bmatrix}
\end{align*}
%
To understand how we build a transformation bridge from $(A,B)$ directly to $(\tilde{A}, \tilde{B})$, we can visualize a common intermediate destination. Instead of guessing the direct transformation $T$, we route both systems through the \textbf{Intermediate Companion Form} $(\bar{A}, \bar{B})$:

\vspace{0.15in}
\begin{equation*}
\begin{array}{ccc}
\text{\textbf{Original System}} & & \text{\textbf{Reachable Canonical Form}} \\[1ex]
(A, B) & \xrightarrow{\qquad \qquad T \ = \ \mathbf{R} \tilde{\mathbf{R}}^{-1} \qquad \qquad} & (\tilde{A}, \tilde{B}) \\[2ex]
\mathbf{R} \searrow & & \swarrow \tilde{\mathbf{R}} \\[2ex]
& \text{\textbf{Intermediate Companion Form}} & \\
& (\bar{A}, \bar{B}) &
\end{array}
\end{equation*}
\vspace{0.15in}

Here is how the logic follows the diagram:
\begin{enumerate}
    \item We established that applying the reachability matrix $\mathbf{R}$ to the original system $(A,B)$ brings it to the intermediate form $(\bar{A}, \bar{B})$. 
    \begin{align*}
        \bar{A} = \mathbf{R}^{-1} A \mathbf{R} \quad \text{and} \quad \bar{B} = \mathbf{R}^{-1} B 
    \end{align*}
    \item Now, what happens if we take the reachable canonical form $(\tilde{A}, \tilde{B})$ and apply \textit{its own} reachability matrix $\tilde{\mathbf{R}}$ as a transformation? Because of the strictly shifting structure of canonical forms, it naturally results in the exact same intermediate companion form! 
    \begin{align*}
        \bar{A} = \tilde{\mathbf{R}}^{-1} \tilde{A} \tilde{\mathbf{R}} \quad \text{and} \quad \bar{B} = \tilde{\mathbf{R}}^{-1} \tilde{B} 
    \end{align*}
    \item By equating the two paths meeting at $(\bar{A}, \bar{B})$, we successfully bridge the gap to find the direct similarity transformation $T$:
    \begin{align*}
        \mathbf{R}^{-1} A \mathbf{R} = \tilde{\mathbf{R}}^{-1} \tilde{A} \tilde{\mathbf{R}} \implies \tilde{A} = (\mathbf{R} \tilde{\mathbf{R}}^{-1})^{-1} A (\mathbf{R} \tilde{\mathbf{R}}^{-1})
    \end{align*}
\end{enumerate}

Thus, if we define the overall coordinate transformation matrix as $T_{\mathbf{R} \tilde{\mathbf{R}}^{-1}} = \mathbf{R} \tilde{\mathbf{R}}^{-1}$ (mapping the original system directly to the reachable canonical form), we obtain the new coordinates:
%
\begin{align*}
	q &= \left( T_{\mathbf{R} \tilde{\mathbf{R}}^{-1}} \right)^{-1} x \implies \dot{q} = \tilde{A} q + \tilde{B} u \\
	\tilde{A} &= \left( T_{\mathbf{R} \tilde{\mathbf{R}}^{-1}} \right)^{-1} \, A \, T_{\mathbf{R} \tilde{\mathbf{R}}^{-1}} \ , \ \tilde{B} = \left(  T_{\mathbf{R} \tilde{\mathbf{R}}^{-1}} \right)^{-1} B 
\end{align*}
%
Now let's apply the state feedback rule in these transformed coordinates $u = \gamma r - \tilde{K} q$:
%
\begin{align*}
 \dot{q} &= ( \tilde{A} - \tilde{B} \tilde{K} ) q + \gamma \tilde{B} r 
\end{align*}
Because of the elegant structure of $\tilde{B}$, the product $\tilde{B} \tilde{K}$ only affects the \textit{last row} of $\tilde{A}$:
\begin{align*}
 ( \tilde{A} - \tilde{B} \tilde{K} ) &= 
 \left[ \begin{array}{c|c|c|c|c|c} 
0 & 1 & 0 & 0 & \cdots & 0 \\
0 & 0 & 1 &  0 & \cdots & 0 \\
\vdots & \vdots & \vdots & \ddots &  & \vdots \\
0 & 0 & 0 & \cdots  & 1 & 0 \\
0 & 0 & 0 & \cdots  & 0 & 1 \\
-a_{0} & -a_{1} & -a_{2} & \cdots  & -a_{n-2} & -a_{n-1} 
\end{array} \right] - \begin{bmatrix} 0 \\ 0 \\ \vdots \\ 0 \\ 1 \end{bmatrix}
\begin{bmatrix} \tilde{k}_0 & \tilde{k}_1 & \cdots & \tilde{k}_{n-1} \end{bmatrix}
\\
&=  \left[ \begin{array}{c|c|c|c|c|c} 
0 & 1 & 0 & 0 & \cdots & 0 \\
0 & 0 & 1 &  0 & \cdots & 0 \\
\vdots & \vdots & \vdots & \ddots &  & \vdots \\
0 & 0 & 0 & \cdots  & 1 & 0 \\
0 & 0 & 0 & \cdots  & 0 & 1 \\
-(a_{0}+\tilde{k}_0) & -(a_{1}+\tilde{k}_1) & -(a_{2}+\tilde{k}_2) & \cdots  & -(a_{n-2}+\tilde{k}_{n-2}) & -(a_{n-1} + \tilde{k}_{n-1})
\end{array} \right]
\end{align*}
%
Reading the coefficients directly from the last row, the closed-loop characteristic equation becomes:
\begin{align*}
\mathrm{det}\left( \lambda I - (\tilde{A} - \tilde{B} \tilde{K}) \right) &=
\lambda^n + ( a_{n-1} + \tilde{k}_{n-1} ) \lambda^{n-1}
	+ \cdots + ( a_{1} + \tilde{k}_1 )  \lambda + ( a_0 + \tilde{k}_0 )
\end{align*}
%
It is now glaringly evident that to match the desired polynomial coefficients $a_i^*$, we simply choose $\tilde{k}_i$:
%
\begin{align*}
	\tilde{k}_i = a_i^* - a_i \quad \forall i \in \{0, \dots, n-1\}
\end{align*} 
%
Finally, we map this feedback gain back to the original coordinate frame:
\begin{align*}
 u &= \gamma r - \tilde{K} q = \gamma r - \tilde{K} \left( T_{\mathbf{R} \tilde{\mathbf{R}}^{-1}} \right)^{-1} x  \implies K = \tilde{K} \left( T_{\mathbf{R} \tilde{\mathbf{R}}^{-1}} \right)^{-1}
\end{align*} 
%
Since similarity transformations preserve eigenvalues, $(A - B K)$ and $( \tilde{A} - \tilde{B} \tilde{K} )$ share the exact same characteristic polynomial. This completes the proof for the SISO case.

\vspace{0.2in}
\textbf{Key Insight - Connection to Ackermann's Formula:} 
The derivation above provides an explicit algorithm to calculate $K$. If we substitute the algebra, this exact process simplifies gracefully into \textbf{Ackermann's Formula}, a famous direct computation method for pole placement in SISO systems:
\begin{align*}
    K = \begin{bmatrix} 0 & 0 & \dots & 0 & 1 \end{bmatrix} \mathcal{R}^{-1} \phi_d(A)
\end{align*}
where $\mathcal{R}$ is the reachability matrix and $\phi_d(A)$ is the desired closed-loop characteristic polynomial evaluated at the matrix $A$.

 \section{Stabilizability}

 \textbf{Intuition \& Insight:} Reachability is a very strong condition; it demands total authority to drive the system to \textit{any} arbitrary state. \textbf{Stabilizability} is a more practical, relaxed condition. It recognizes that if we lack control over certain parts of the state space, the system is still viable \textit{if and only if} those uncontrollable parts naturally decay to the origin on their own.

 \textbf{Definition:} An LTI system is called stabilizable if 
 $\forall x_0 \in \mathbb{R}^{n}$, $\exists \, u(t)$ (or $u[k]$) such that
 $\lim_{t\to\infty} x(t) = 0$ (or $\lim_{k\to\infty} x[k] = 0$).

 An obvious \textbf{sufficient} condition on stabilizability for CT-LTI and DT-LTI systems is that if $(A,B)$ is reachable, then $(A,B)$ is stabilizable, since
 in such a case $\forall x_0 \in \mathbb{R}^{n}$, $\exists \, u(t)$ or $u[k]$
 such that $x(T) = 0 \, \forall T > 0$ for CT systems, and $x[k] = 0 \, \forall  k \geq n$ for DT systems. 

Another \textbf{sufficient} (but weaker) condition on stabilizability for DT-LTI systems is that if $(A,B)$ is controllable (can reach the origin, even if it can't reach everywhere else), then $(A,B)$ is stabilizable. 

\textbf{Theorem:} The following statements are rigorously equivalent:
\begin{itemize}
\item $(A,B)$ is stabilizable $\iff$
\item Unreachable modes of $(A,B)$ are strictly asymptotically stable $\iff$
\item \textbf{PBH Eigenvector Test:} $\forall (\lambda_i , w_i)$ such that $w_i^T A = \lambda_i w_i^T$
and $\mathrm{Re}\lbrace \lambda_i \rbrace \geq 0$ (or $|\lambda_i| \geq 1$ for DT), we must have $w_i^T B \neq 0$ $\iff$
\item \textbf{PBH Rank Test:} $\mathrm{rank}\left[ A - \lambda I \, | \, B \right] = n$,
$\forall \lambda \in \mathbb{C}$ such that $\mathrm{Re}\lbrace \lambda \rbrace \geq 0$ for CT systems (or $\forall \lambda \in \mathbb{C}$ such that $|\lambda|  \geq 1$ for DT systems).
\end{itemize}

\textbf{Proof:} Given $(A,B)$ that is not necessarily reachable, we can utilize the Kalman standard unreachable form to isolate the reachable and unreachable dynamics:
%
\begin{align*}
\dot{x} = \left[ \begin{array}{c} \dot{x}_r \\ \dot{x}_{\bar{r}} \end{array} \right]  = \left[ \begin{array}{c|c} A_{r} & A_{12} \\ \hline 0 & A_{\bar{r}} \end{array} \right] \left[ \begin{array}{c} x_r \\ x_{\bar{r}} \end{array} \right]  
+ \left[ \begin{array}{c} B_r \\ 0 \end{array} \right] u
\end{align*}
%
If the unforced subsystem $\dot{x}_{\bar{r}} = A_{\bar{r}} x_{\bar{r}}$ is \textit{not} asymptotically stable, then obviously the whole system is unstable regardless of $u(t)$ because the input has no pathway to influence $x_{\bar{r}}$; thus, the system is not stabilizable.
%
Now, assume $\dot{x}_{\bar{r}} = A_{\bar{r}} x_{\bar{r}}$ \textit{is} asymptotically stable. Let us apply a state feedback control relying only on the reachable states, $u = - K_{r} x_{r}$. The closed-loop system takes the form:
%
\begin{align*}
\dot{x} = \left[ \begin{array}{c} \dot{x}_r \\ \dot{x}_{\bar{r}} \end{array} \right]  = \underbrace{\left[ \begin{array}{c|c} A_{r} - B_r K_r & A_{12} \\ \hline 0 & A_{\bar{r}} \end{array} \right]}_{A_{cl}} \left[ \begin{array}{c} x_r \\ x_{\bar{r}} \end{array} \right]  
\end{align*}
%
Because $A_{cl}$ is block upper-triangular, its closed-loop spectrum (eigenvalues) is simply the union of the block diagonal spectra:
\begin{align*}
    \sigma(A_{cl}) = \sigma(A_{r} - B_r K_r) \cup \sigma(A_{\bar{r}})
\end{align*}
Since the pair $(A_r,B_r)$ is fully reachable by definition, we can find a gain $K_r$ such that the eigenvalues of $(A_{r} - B_r K_r)$ are located deep in the open left-half of the complex plane. Since $\sigma(A_{\bar{r}})$ is already stable, the union $\sigma(A_{cl})$ is entirely stable. Thus, the system is stabilizable. The exact same logic applies to DT systems.

% **** This ENDS THE EXAMPLES. DON'T DELETE THE FOLLOWING LINE:
\end{document}